\section{Legal increasing appends preserve every existing first positive slope}
For a zero-extended ordinary rank sequence write
$\delta F_n=F_n-F_{n-1}$ and $D(F)_n=(\delta F_n)^+$.
The harmless initial coefficient $D(F)_0=1$ may be included;
the usual convention $D_0=D_1=0$ gives the same theorem.

\begin{theorem}[Coefficientwise first-derivative order]
For every increasing normalized arm $X=1^\alpha U$ and every
legal append $t\ge\max U$, $t\ge2$,
\[
                         D(F_{Xt})_n\ge D(F_X)_n
                         \qquad\hbox{for every }n.       \tag{1}
\]
Moreover, the added-word polynomial $F_{Xt}-F_X$ has
nonnegative first differences throughout the old rising prefix.
The statement is unbounded in the number and sizes of all parts.
\end{theorem}
\begin{proof}
The all-ones base has no positive slope after rank zero, so the
first non-one append satisfies (1). For the induction step write
$X=ZL$ and assume the theorem for the preceding append.
The exact TWO-stage identity is
\[
 F_{Xt}-F_X=S_L(F_X-F_Z)+(S_t-S_L)F_X.                  \tag{2}
\]
Both kernels are nonnegative and supported in positive degrees:
$S_t-S_L=z^{L+1}[t-L]$. First let $n$ belong to the old
nondecreasing prefix, including its plateaus. Then
$\delta F_X(j)\ge0$ for every
$j\le n$. At every such $j$, the induction hypothesis gives
$\delta F_X(j)-\delta F_Z(j)\ge0$: if the latter old slope
is positive use (1), and otherwise use $\delta F_X(j)\ge0$.
Apply $1-z$ to (2). Every coefficient sampled on its right side
at rank $n$ is one of these nonnegative earlier coefficients.
Thus $\delta F_{Xt}(n)-\delta F_X(n)\ge0$. Outside that
prefix the old slope is nonpositive by unimodality, so (1)
is immediate from positive clipping. This
proves the induction, including the initial and terminal ranks.
\end{proof}

\paragraph{The second positive derivative is not ordered.}
The analogous assertion for $W=D(D(F))$ is false, even within
the already solved small-peak original families. For every
$\alpha,h\ge1$ put
\[
 X=1^\alpha 2^h\,4\,5^4,\qquad X'=X\,6.
\]
The coefficients through rank seven are, for ALL these parameters,
\[
 F_X=(1,1,2,4,8,16,31,61,\ldots),\qquad
 F_{X'}=(1,1,2,4,8,16,32,62,\ldots).
\]
Hence $D_X(6,7)=(15,30)$ and $D_{X'}(6,7)=(16,30)$,
so $W_X(7)=15>14=W_{X'}(7)$.
More generally one may replace $1^\alpha2^h$ by ANY nonrectangular
arm $Z$ with maximum at most four. Only its initial coefficients
$(1,1,2)$ can affect these displayed ranks: the five suffix
appends shift every dependence on $F_Z$ by at least five.

The first-order theorem does not imply likelihood-ratio order,
second-order order, C, S, or original-core unimodality. In particular
the exact positive TWO-stage source in (2) must be retained when
attempting any stronger propagation argument.
